% --- Archon physics formalization source begin ---
% archon:physics
% archon:covers PhyXMiniProblems/problem_phyx_mini_0326.lean
% archon:source-report reports/phyx_mini/problem_phyx_mini_0326.source.json

\chapter{Physics problem phyx\_mini\_0326}
\label{ch:PhyXMiniProblems_problem_phyx_mini_0326}

\paragraph{Problem source.}
\#\# Physical scenario

On July 10, 1996, a granite block broke away from a wall in Yosemite Valley and, as it began to slide down the wall, was launched into projectile motion. Seismic waves produced by its impact with the ground triggered seismographs as far away as \$200\$ km. Later measurements indicated that the block had a mass between \$7.3 \textbackslash{}times 10\textasciicircum{}7\$ kg and \$1.7 \textbackslash{}times 10\textasciicircum{}8\$ kg and that it landed \$500\$ m vertically below the launch point and \$30\$ m horizontally from it. Consider two types of seismic waves that spread from the impact point—a hemispherical \textbackslash{}textit\{body wave\} traveled through the ground in an expanding hemisphere and a cylindrical \textbackslash{}textit\{surface wave\} traveled along the ground in an expanding shallow vertical cylinder. Assume that the impact lasted \textbackslash{}(0.50\textbackslash{},\textbackslash{}mathrm\{s\}\textbackslash{}), the vertical cylinder had a depth \textbackslash{}(d = 5.0\textbackslash{},\textbackslash{}mathrm\{m\}\textbackslash{}), and each wave type received 20\textbackslash{}\% of the energy the block had just before impact.

\#\# Question

Neglecting any mechanical energy loss the waves experienced as they traveled, determine the intensities of the surface wave when they reached a seismograph \textbackslash{}(200\textbackslash{},\textbackslash{}mathrm\{km\}\textbackslash{}) away.

\#\# Answer choices

- A: A: \$43\$ kW/m\$\textasciicircum{}2\$

- B: B: \$48\$ kW/m\$\textasciicircum{}2\$

- C: C: \$53\$ kW/m\$\textasciicircum{}2\$

- D: D: \$58\$ kW/m\$\textasciicircum{}2\$

\#\# Figure caption (auxiliary; use the image as primary evidence)

The figure illustrates a diagram depicting wave propagation in a medium.

- **Cylindrical Wave**: Shown as a red rectangular line, indicating wave propagation horizontally in both directions. Arrows point outward from the source of impact.
- **Hemispherical Wave**: Depicted as a blue semicircular line, indicating wave propagation downward and outward in a semicircular pattern. Arrows also point outward.
- **Impact**: Represented by a black dot on the surface to indicate the origin of the waves.
- **Depth**: Labeled as "d" with a double-headed vertical arrow on the right side, showing the distance from the surface to deeper regions of the medium.
- The surface is illustrated as a horizontal line above the shaded area, suggesting the medium's material through which the waves are propagating.

The relationships shown include the wave propagation directions and the origin point of the impact.

\#\# Dataset metadata

Wave Properties; Physical Model Grounding Reasoning, Multi-Formula Reasoning

\paragraph{Recorded answer/context.}
D: D: \$58\$ kW/m\$\textasciicircum{}2\$

\paragraph{Figure/image path.}
/root/proposal\_for\_physic/science-mango/phyx\_mini\_run/phyx\_data/test\_image/326.png

\paragraph{Formalization target.}
create a compiling Lean file with sorry bodies at `PhyXMiniProblems/problem\_phyx\_mini\_0326.lean`. The Lean declarations must preserve the physical quantities, dimensions or dimensional roles, figure labels, governing-law hypotheses, and final relation expressed by this problem.
Use Mathlib/Physlib names found through LeanExplore where available. If a physics API is missing, introduce faithful local abstractions rather than scalar placeholder aliases.

\begin{theorem}[Physics formalization target]
\label{thm:physics:phyx_mini_0326:target}
\lean{PhyXMiniProblems.ProblemPhyXMini0326.problem_phyx_mini_0326}
\uses{def:physics:phyx-mini-0326:phyxminiproblems-problemphyxmini0326-lengthinmeters, def:physics:phyx-mini-0326:phyxminiproblems-problemphyxmini0326-timeinseconds, def:physics:phyx-mini-0326:phyxminiproblems-problemphyxmini0326-massinkilograms, def:physics:phyx-mini-0326:phyxminiproblems-problemphyxmini0326-accelerationinmeterspersecondsquared, def:physics:phyx-mini-0326:phyxminiproblems-problemphyxmini0326-intensityinkilowattspersquaremeter, def:physics:phyx-mini-0326:phyxminiproblems-problemphyxmini0326-yosemiterockfallsetup, def:physics:phyx-mini-0326:phyxminiproblems-problemphyxmini0326-hasphysicalrockfallparameters, def:physics:phyx-mini-0326:phyxminiproblems-problemphyxmini0326-matchesyosemiterockfallstatementandfigure, def:physics:phyx-mini-0326:phyxminiproblems-problemphyxmini0326-satisfiesrockfallimpactenergylaw, def:physics:phyx-mini-0326:phyxminiproblems-problemphyxmini0326-satisfiesseismicenergypartitionlaw, def:physics:phyx-mini-0326:phyxminiproblems-problemphyxmini0326-satisfiesfiniteimpactpowerlaw, def:physics:phyx-mini-0326:phyxminiproblems-problemphyxmini0326-satisfiesseismicwavefrontgeometrylaw, def:physics:phyx-mini-0326:phyxminiproblems-problemphyxmini0326-satisfieslosslessseismicintensitylaw, def:physics:phyx-mini-0326:phyxminiproblems-problemphyxmini0326-displayedintensityinkilowattspersquaremeter, def:physics:phyx-mini-0326:phyxminiproblems-problemphyxmini0326-recordeddatasetanswerchoice, def:physics:phyx-mini-0326:phyxminiproblems-problemphyxmini0326-roundstodisplayedwholekilowattpersquaremeter}
The assigned autoformalize agent should translate this physics problem into Lean declarations in the covered file, with theorem and lemma proof bodies written as `by sorry`.
\end{theorem}
\begin{proof}
This is an autoformalization task, not a proof task. Produce statements that can later be proved by the physics prover without weakening the physical model.
\end{proof}

% --- Archon declaration topology begin ---
\section{Declaration topology}
\begin{definition}[Mass Quantity]
  \label{def:physics:phyx-mini-0326:phyxminiproblems-problemphyxmini0326-massquantity}
  \lean{PhyXMiniProblems.ProblemPhyXMini0326.MassQuantity}
  A nonnegative physical mass
\end{definition}
\begin{proof}
  This is the declared model definition of Mass Quantity.
\end{proof}
\begin{definition}[Length Quantity]
  \label{def:physics:phyx-mini-0326:phyxminiproblems-problemphyxmini0326-lengthquantity}
  \lean{PhyXMiniProblems.ProblemPhyXMini0326.LengthQuantity}
  A nonnegative physical length
\end{definition}
\begin{proof}
  This is the declared model definition of Length Quantity.
\end{proof}
\begin{definition}[Time Quantity]
  \label{def:physics:phyx-mini-0326:phyxminiproblems-problemphyxmini0326-timequantity}
  \lean{PhyXMiniProblems.ProblemPhyXMini0326.TimeQuantity}
  A nonnegative physical duration
\end{definition}
\begin{proof}
  This is the declared model definition of Time Quantity.
\end{proof}
\begin{definition}[Acceleration Quantity]
  \label{def:physics:phyx-mini-0326:phyxminiproblems-problemphyxmini0326-accelerationquantity}
  \lean{PhyXMiniProblems.ProblemPhyXMini0326.AccelerationQuantity}
  A nonnegative acceleration magnitude
\end{definition}
\begin{proof}
  This is the declared model definition of Acceleration Quantity.
\end{proof}
\begin{definition}[Energy Quantity]
  \label{def:physics:phyx-mini-0326:phyxminiproblems-problemphyxmini0326-energyquantity}
  \lean{PhyXMiniProblems.ProblemPhyXMini0326.EnergyQuantity}
  Mechanical or seismic energy, using Physlib's dimensional energy type
\end{definition}
\begin{proof}
  This is the declared model definition of Energy Quantity.
\end{proof}
\begin{definition}[Area Quantity]
  \label{def:physics:phyx-mini-0326:phyxminiproblems-problemphyxmini0326-areaquantity}
  \lean{PhyXMiniProblems.ProblemPhyXMini0326.AreaQuantity}
  Wavefront area, using Physlib's dimensional area type
\end{definition}
\begin{proof}
  This is the declared model definition of Area Quantity.
\end{proof}
\begin{definition}[Power Quantity]
  \label{def:physics:phyx-mini-0326:phyxminiproblems-problemphyxmini0326-powerquantity}
  \lean{PhyXMiniProblems.ProblemPhyXMini0326.PowerQuantity}
  A nonnegative power, with coherent SI unit watt
\end{definition}
\begin{proof}
  This is the declared model definition of Power Quantity.
\end{proof}
\begin{definition}[Intensity Quantity]
  \label{def:physics:phyx-mini-0326:phyxminiproblems-problemphyxmini0326-intensityquantity}
  \lean{PhyXMiniProblems.ProblemPhyXMini0326.IntensityQuantity}
  Intensity is power per area, so its physical dimension is (M L² T⁻³) / L² = M T⁻³; its coherent SI unit is watt per square metre
\end{definition}
\begin{proof}
  This is the declared model definition of Intensity Quantity.
\end{proof}
\begin{definition}[nonnegative Quantity Readout]
  \label{def:physics:phyx-mini-0326:phyxminiproblems-problemphyxmini0326-nonnegativequantityreadout}
  \lean{PhyXMiniProblems.ProblemPhyXMini0326.nonnegativeQuantityReadout}
  Read a nonnegative dimensional quantity in a coherent unit system
\end{definition}
\begin{proof}
  This is the declared model definition of nonnegative Quantity Readout.
\end{proof}
\begin{definition}[length Readout]
  \label{def:physics:phyx-mini-0326:phyxminiproblems-problemphyxmini0326-lengthreadout}
  \lean{PhyXMiniProblems.ProblemPhyXMini0326.lengthReadout}
  \uses{def:physics:phyx-mini-0326:phyxminiproblems-problemphyxmini0326-lengthquantity, def:physics:phyx-mini-0326:phyxminiproblems-problemphyxmini0326-nonnegativequantityreadout}
  Read a physical length in a selected length unit
\end{definition}
\begin{proof}
  This is the declared model definition of length Readout.
\end{proof}
\begin{definition}[length In Meters]
  \label{def:physics:phyx-mini-0326:phyxminiproblems-problemphyxmini0326-lengthinmeters}
  \lean{PhyXMiniProblems.ProblemPhyXMini0326.lengthInMeters}
  \uses{def:physics:phyx-mini-0326:phyxminiproblems-problemphyxmini0326-lengthquantity, def:physics:phyx-mini-0326:phyxminiproblems-problemphyxmini0326-lengthreadout}
  Metre readout of a physical length
\end{definition}
\begin{proof}
  This is the declared model definition of length In Meters.
\end{proof}
\begin{definition}[length In Kilometers]
  \label{def:physics:phyx-mini-0326:phyxminiproblems-problemphyxmini0326-lengthinkilometers}
  \lean{PhyXMiniProblems.ProblemPhyXMini0326.lengthInKilometers}
  \uses{def:physics:phyx-mini-0326:phyxminiproblems-problemphyxmini0326-lengthquantity, def:physics:phyx-mini-0326:phyxminiproblems-problemphyxmini0326-lengthreadout}
  Kilometre readout of a physical length
\end{definition}
\begin{proof}
  This is the declared model definition of length In Kilometers.
\end{proof}
\begin{definition}[time In Seconds]
  \label{def:physics:phyx-mini-0326:phyxminiproblems-problemphyxmini0326-timeinseconds}
  \lean{PhyXMiniProblems.ProblemPhyXMini0326.timeInSeconds}
  \uses{def:physics:phyx-mini-0326:phyxminiproblems-problemphyxmini0326-timequantity, def:physics:phyx-mini-0326:phyxminiproblems-problemphyxmini0326-nonnegativequantityreadout}
  Second readout of a physical duration
\end{definition}
\begin{proof}
  This is the declared model definition of time In Seconds.
\end{proof}
\begin{definition}[mass In Kilograms]
  \label{def:physics:phyx-mini-0326:phyxminiproblems-problemphyxmini0326-massinkilograms}
  \lean{PhyXMiniProblems.ProblemPhyXMini0326.massInKilograms}
  \uses{def:physics:phyx-mini-0326:phyxminiproblems-problemphyxmini0326-massquantity, def:physics:phyx-mini-0326:phyxminiproblems-problemphyxmini0326-nonnegativequantityreadout}
  Kilogram readout of a physical mass
\end{definition}
\begin{proof}
  This is the declared model definition of mass In Kilograms.
\end{proof}
\begin{definition}[acceleration In Meters Per Second Squared]
  \label{def:physics:phyx-mini-0326:phyxminiproblems-problemphyxmini0326-accelerationinmeterspersecondsquared}
  \lean{PhyXMiniProblems.ProblemPhyXMini0326.accelerationInMetersPerSecondSquared}
  \uses{def:physics:phyx-mini-0326:phyxminiproblems-problemphyxmini0326-accelerationquantity, def:physics:phyx-mini-0326:phyxminiproblems-problemphyxmini0326-nonnegativequantityreadout}
  Metre-per-second-squared readout of an acceleration
\end{definition}
\begin{proof}
  This is the declared model definition of acceleration In Meters Per Second Squared.
\end{proof}
\begin{definition}[energy In Joules]
  \label{def:physics:phyx-mini-0326:phyxminiproblems-problemphyxmini0326-energyinjoules}
  \lean{PhyXMiniProblems.ProblemPhyXMini0326.energyInJoules}
  \uses{def:physics:phyx-mini-0326:phyxminiproblems-problemphyxmini0326-energyquantity}
  Joule readout of an energy
\end{definition}
\begin{proof}
  This is the declared model definition of energy In Joules.
\end{proof}
\begin{definition}[area In Square Meters]
  \label{def:physics:phyx-mini-0326:phyxminiproblems-problemphyxmini0326-areainsquaremeters}
  \lean{PhyXMiniProblems.ProblemPhyXMini0326.areaInSquareMeters}
  \uses{def:physics:phyx-mini-0326:phyxminiproblems-problemphyxmini0326-areaquantity, def:physics:phyx-mini-0326:phyxminiproblems-problemphyxmini0326-nonnegativequantityreadout}
  Square-metre readout of a wavefront area
\end{definition}
\begin{proof}
  This is the declared model definition of area In Square Meters.
\end{proof}
\begin{definition}[power In Watts]
  \label{def:physics:phyx-mini-0326:phyxminiproblems-problemphyxmini0326-powerinwatts}
  \lean{PhyXMiniProblems.ProblemPhyXMini0326.powerInWatts}
  \uses{def:physics:phyx-mini-0326:phyxminiproblems-problemphyxmini0326-powerquantity, def:physics:phyx-mini-0326:phyxminiproblems-problemphyxmini0326-nonnegativequantityreadout}
  Watt readout of a wave power
\end{definition}
\begin{proof}
  This is the declared model definition of power In Watts.
\end{proof}
\begin{definition}[intensity In Watts Per Square Meter]
  \label{def:physics:phyx-mini-0326:phyxminiproblems-problemphyxmini0326-intensityinwattspersquaremeter}
  \lean{PhyXMiniProblems.ProblemPhyXMini0326.intensityInWattsPerSquareMeter}
  \uses{def:physics:phyx-mini-0326:phyxminiproblems-problemphyxmini0326-intensityquantity, def:physics:phyx-mini-0326:phyxminiproblems-problemphyxmini0326-nonnegativequantityreadout}
  Watt-per-square-metre readout of a wave intensity
\end{definition}
\begin{proof}
  This is the declared model definition of intensity In Watts Per Square Meter.
\end{proof}
\begin{definition}[intensity In Kilowatts Per Square Meter]
  \label{def:physics:phyx-mini-0326:phyxminiproblems-problemphyxmini0326-intensityinkilowattspersquaremeter}
  \lean{PhyXMiniProblems.ProblemPhyXMini0326.intensityInKilowattsPerSquareMeter}
  \uses{def:physics:phyx-mini-0326:phyxminiproblems-problemphyxmini0326-intensityquantity, def:physics:phyx-mini-0326:phyxminiproblems-problemphyxmini0326-intensityinwattspersquaremeter}
  Kilowatt-per-square-metre readout used by the answer choices
\end{definition}
\begin{proof}
  This is the declared model definition of intensity In Kilowatts Per Square Meter.
\end{proof}
\begin{definition}[Rockfall Event]
  \label{def:physics:phyx-mini-0326:phyxminiproblems-problemphyxmini0326-rockfallevent}
  \lean{PhyXMiniProblems.ProblemPhyXMini0326.RockfallEvent}
  The historical event named in the problem statement
\end{definition}
\begin{proof}
  This is the declared model definition of Rockfall Event.
\end{proof}
\begin{definition}[Block Material]
  \label{def:physics:phyx-mini-0326:phyxminiproblems-problemphyxmini0326-blockmaterial}
  \lean{PhyXMiniProblems.ProblemPhyXMini0326.BlockMaterial}
  Material of the falling block
\end{definition}
\begin{proof}
  This is the declared model definition of Block Material.
\end{proof}
\begin{definition}[Seismic Wave Kind]
  \label{def:physics:phyx-mini-0326:phyxminiproblems-problemphyxmini0326-seismicwavekind}
  \lean{PhyXMiniProblems.ProblemPhyXMini0326.SeismicWaveKind}
  The two seismic-wave species produced by the impact
\end{definition}
\begin{proof}
  This is the declared model definition of Seismic Wave Kind.
\end{proof}
\begin{definition}[Wavefront Geometry]
  \label{def:physics:phyx-mini-0326:phyxminiproblems-problemphyxmini0326-wavefrontgeometry}
  \lean{PhyXMiniProblems.ProblemPhyXMini0326.WavefrontGeometry}
  Idealized wavefront shapes stated in the prose and shown in the image
\end{definition}
\begin{proof}
  This is the declared model definition of Wavefront Geometry.
\end{proof}
\begin{definition}[Figure Propagation Direction]
  \label{def:physics:phyx-mini-0326:phyxminiproblems-problemphyxmini0326-figurepropagationdirection}
  \lean{PhyXMiniProblems.ProblemPhyXMini0326.FigurePropagationDirection}
  Qualitative propagation directions displayed by the arrows in the image
\end{definition}
\begin{proof}
  This is the declared model definition of Figure Propagation Direction.
\end{proof}
\begin{definition}[Figure Feature]
  \label{def:physics:phyx-mini-0326:phyxminiproblems-problemphyxmini0326-figurefeature}
  \lean{PhyXMiniProblems.ProblemPhyXMini0326.FigureFeature}
  Named locations or markers visible in the primary image
\end{definition}
\begin{proof}
  This is the declared model definition of Figure Feature.
\end{proof}
\begin{definition}[Seismic Propagation Figure]
  \label{def:physics:phyx-mini-0326:phyxminiproblems-problemphyxmini0326-seismicpropagationfigure}
  \lean{PhyXMiniProblems.ProblemPhyXMini0326.SeismicPropagationFigure}
  \uses{def:physics:phyx-mini-0326:phyxminiproblems-problemphyxmini0326-lengthquantity, def:physics:phyx-mini-0326:phyxminiproblems-problemphyxmini0326-seismicwavekind, def:physics:phyx-mini-0326:phyxminiproblems-problemphyxmini0326-wavefrontgeometry, def:physics:phyx-mini-0326:phyxminiproblems-problemphyxmini0326-figurepropagationdirection, def:physics:phyx-mini-0326:phyxminiproblems-problemphyxmini0326-figurefeature}
  Primary-image evidence is stored separately from the physical laws. The figure carries the depth label d but does not itself assign its numerical value; that value comes from the prose data predicate below
\end{definition}
\begin{proof}
  This is the declared model definition of Seismic Propagation Figure.
\end{proof}
\begin{definition}[Yosemite Rockfall Setup]
  \label{def:physics:phyx-mini-0326:phyxminiproblems-problemphyxmini0326-yosemiterockfallsetup}
  \lean{PhyXMiniProblems.ProblemPhyXMini0326.YosemiteRockfallSetup}
  \uses{def:physics:phyx-mini-0326:phyxminiproblems-problemphyxmini0326-massquantity, def:physics:phyx-mini-0326:phyxminiproblems-problemphyxmini0326-lengthquantity, def:physics:phyx-mini-0326:phyxminiproblems-problemphyxmini0326-timequantity, def:physics:phyx-mini-0326:phyxminiproblems-problemphyxmini0326-accelerationquantity, def:physics:phyx-mini-0326:phyxminiproblems-problemphyxmini0326-energyquantity, def:physics:phyx-mini-0326:phyxminiproblems-problemphyxmini0326-areaquantity, def:physics:phyx-mini-0326:phyxminiproblems-problemphyxmini0326-powerquantity, def:physics:phyx-mini-0326:phyxminiproblems-problemphyxmini0326-intensityquantity, def:physics:phyx-mini-0326:phyxminiproblems-problemphyxmini0326-rockfallevent, def:physics:phyx-mini-0326:phyxminiproblems-problemphyxmini0326-blockmaterial, def:physics:phyx-mini-0326:phyxminiproblems-problemphyxmini0326-seismicwavekind, def:physics:phyx-mini-0326:phyxminiproblems-problemphyxmini0326-seismicpropagationfigure}
  The independent quantities and observables in the rockfall model. The actual block mass is distinct from the two measured endpoints. Wave energy, power, area, and intensity are also stored independently, then related by governing laws rather than by answer-producing definitions
\end{definition}
\begin{proof}
  This is the declared model definition of Yosemite Rockfall Setup.
\end{proof}
\begin{definition}[Has Physical Rockfall Parameters]
  \label{def:physics:phyx-mini-0326:phyxminiproblems-problemphyxmini0326-hasphysicalrockfallparameters}
  \lean{PhyXMiniProblems.ProblemPhyXMini0326.HasPhysicalRockfallParameters}
  \uses{def:physics:phyx-mini-0326:phyxminiproblems-problemphyxmini0326-lengthinmeters, def:physics:phyx-mini-0326:phyxminiproblems-problemphyxmini0326-timeinseconds, def:physics:phyx-mini-0326:phyxminiproblems-problemphyxmini0326-massinkilograms, def:physics:phyx-mini-0326:phyxminiproblems-problemphyxmini0326-accelerationinmeterspersecondsquared, def:physics:phyx-mini-0326:phyxminiproblems-problemphyxmini0326-energyinjoules, def:physics:phyx-mini-0326:phyxminiproblems-problemphyxmini0326-powerinwatts, def:physics:phyx-mini-0326:phyxminiproblems-problemphyxmini0326-yosemiterockfallsetup}
  Positivity and nondegeneracy of the independent physical parameters
\end{definition}
\begin{proof}
  This is the declared model definition of Has Physical Rockfall Parameters.
\end{proof}
\begin{definition}[Matches Yosemite Rockfall Statement And Figure]
  \label{def:physics:phyx-mini-0326:phyxminiproblems-problemphyxmini0326-matchesyosemiterockfallstatementandfigure}
  \lean{PhyXMiniProblems.ProblemPhyXMini0326.MatchesYosemiteRockfallStatementAndFigure}
  \uses{def:physics:phyx-mini-0326:phyxminiproblems-problemphyxmini0326-lengthinmeters, def:physics:phyx-mini-0326:phyxminiproblems-problemphyxmini0326-lengthinkilometers, def:physics:phyx-mini-0326:phyxminiproblems-problemphyxmini0326-timeinseconds, def:physics:phyx-mini-0326:phyxminiproblems-problemphyxmini0326-massinkilograms, def:physics:phyx-mini-0326:phyxminiproblems-problemphyxmini0326-accelerationinmeterspersecondsquared, def:physics:phyx-mini-0326:phyxminiproblems-problemphyxmini0326-yosemiterockfallsetup}
  Textual measurements and qualitative data, together with the labels and geometry read from the primary image. The actual mass is constrained by the reported interval but is not replaced by either endpoint
\end{definition}
\begin{proof}
  This is the declared model definition of Matches Yosemite Rockfall Statement And Figure.
\end{proof}
\begin{definition}[Satisfies Rockfall Impact Energy Law]
  \label{def:physics:phyx-mini-0326:phyxminiproblems-problemphyxmini0326-satisfiesrockfallimpactenergylaw}
  \lean{PhyXMiniProblems.ProblemPhyXMini0326.SatisfiesRockfallImpactEnergyLaw}
  \uses{def:physics:phyx-mini-0326:phyxminiproblems-problemphyxmini0326-lengthinmeters, def:physics:phyx-mini-0326:phyxminiproblems-problemphyxmini0326-massinkilograms, def:physics:phyx-mini-0326:phyxminiproblems-problemphyxmini0326-accelerationinmeterspersecondsquared, def:physics:phyx-mini-0326:phyxminiproblems-problemphyxmini0326-energyinjoules, def:physics:phyx-mini-0326:phyxminiproblems-problemphyxmini0326-yosemiterockfallsetup}
  Mechanical-energy conservation in the intended rockfall approximation: the impact energy just before collision is the gravitational energy released over the vertical drop. The equation is parameter-generic and contains no wave intensity or answer-choice value
\end{definition}
\begin{proof}
  This is the declared model definition of Satisfies Rockfall Impact Energy Law.
\end{proof}
\begin{definition}[Satisfies Seismic Energy Partition Law]
  \label{def:physics:phyx-mini-0326:phyxminiproblems-problemphyxmini0326-satisfiesseismicenergypartitionlaw}
  \lean{PhyXMiniProblems.ProblemPhyXMini0326.SatisfiesSeismicEnergyPartitionLaw}
  \uses{def:physics:phyx-mini-0326:phyxminiproblems-problemphyxmini0326-energyinjoules, def:physics:phyx-mini-0326:phyxminiproblems-problemphyxmini0326-yosemiterockfallsetup}
  Each wave receives its stated fraction of the pre-impact block energy
\end{definition}
\begin{proof}
  This is the declared model definition of Satisfies Seismic Energy Partition Law.
\end{proof}
\begin{definition}[Satisfies Finite Impact Power Law]
  \label{def:physics:phyx-mini-0326:phyxminiproblems-problemphyxmini0326-satisfiesfiniteimpactpowerlaw}
  \lean{PhyXMiniProblems.ProblemPhyXMini0326.SatisfiesFiniteImpactPowerLaw}
  \uses{def:physics:phyx-mini-0326:phyxminiproblems-problemphyxmini0326-timeinseconds, def:physics:phyx-mini-0326:phyxminiproblems-problemphyxmini0326-energyinjoules, def:physics:phyx-mini-0326:phyxminiproblems-problemphyxmini0326-powerinwatts, def:physics:phyx-mini-0326:phyxminiproblems-problemphyxmini0326-yosemiterockfallsetup}
  Average power is the assigned seismic energy divided by impact duration
\end{definition}
\begin{proof}
  This is the declared model definition of Satisfies Finite Impact Power Law.
\end{proof}
\begin{definition}[Satisfies Seismic Wavefront Geometry Law]
  \label{def:physics:phyx-mini-0326:phyxminiproblems-problemphyxmini0326-satisfiesseismicwavefrontgeometrylaw}
  \lean{PhyXMiniProblems.ProblemPhyXMini0326.SatisfiesSeismicWavefrontGeometryLaw}
  \uses{def:physics:phyx-mini-0326:phyxminiproblems-problemphyxmini0326-lengthinmeters, def:physics:phyx-mini-0326:phyxminiproblems-problemphyxmini0326-areainsquaremeters, def:physics:phyx-mini-0326:phyxminiproblems-problemphyxmini0326-yosemiterockfallsetup}
  Wavefront geometry from the model and figure: the body wave occupies a hemisphere of area 2πr², whereas the shallow surface-wave cylinder has lateral area 2πrd. The latter is the area across which surface-wave power is distributed
\end{definition}
\begin{proof}
  This is the declared model definition of Satisfies Seismic Wavefront Geometry Law.
\end{proof}
\begin{definition}[Satisfies Lossless Seismic Intensity Law]
  \label{def:physics:phyx-mini-0326:phyxminiproblems-problemphyxmini0326-satisfieslosslessseismicintensitylaw}
  \lean{PhyXMiniProblems.ProblemPhyXMini0326.SatisfiesLosslessSeismicIntensityLaw}
  \uses{def:physics:phyx-mini-0326:phyxminiproblems-problemphyxmini0326-lengthinmeters, def:physics:phyx-mini-0326:phyxminiproblems-problemphyxmini0326-areainsquaremeters, def:physics:phyx-mini-0326:phyxminiproblems-problemphyxmini0326-powerinwatts, def:physics:phyx-mini-0326:phyxminiproblems-problemphyxmini0326-intensityinwattspersquaremeter, def:physics:phyx-mini-0326:phyxminiproblems-problemphyxmini0326-yosemiterockfallsetup}
  Lossless propagation preserves the power assigned at impact. Intensity is that transported power divided by the relevant wavefront area. These are general propagation laws at every positive radius, not the requested numerical conclusion at 200 km
\end{definition}
\begin{proof}
  This is the declared model definition of Satisfies Lossless Seismic Intensity Law.
\end{proof}
\begin{lemma}[surface Wave Intensity At Station formula]
  \label{lem:physics:phyx-mini-0326:phyxminiproblems-problemphyxmini0326-surfacewaveintensityatstation-formula}
  \lean{PhyXMiniProblems.ProblemPhyXMini0326.surfaceWaveIntensityAtStation_formula}
  \uses{def:physics:phyx-mini-0326:phyxminiproblems-problemphyxmini0326-lengthinmeters, def:physics:phyx-mini-0326:phyxminiproblems-problemphyxmini0326-timeinseconds, def:physics:phyx-mini-0326:phyxminiproblems-problemphyxmini0326-massinkilograms, def:physics:phyx-mini-0326:phyxminiproblems-problemphyxmini0326-accelerationinmeterspersecondsquared, def:physics:phyx-mini-0326:phyxminiproblems-problemphyxmini0326-intensityinwattspersquaremeter, def:physics:phyx-mini-0326:phyxminiproblems-problemphyxmini0326-yosemiterockfallsetup, def:physics:phyx-mini-0326:phyxminiproblems-problemphyxmini0326-hasphysicalrockfallparameters, def:physics:phyx-mini-0326:phyxminiproblems-problemphyxmini0326-satisfiesrockfallimpactenergylaw, def:physics:phyx-mini-0326:phyxminiproblems-problemphyxmini0326-satisfiesseismicenergypartitionlaw, def:physics:phyx-mini-0326:phyxminiproblems-problemphyxmini0326-satisfiesfiniteimpactpowerlaw, def:physics:phyx-mini-0326:phyxminiproblems-problemphyxmini0326-satisfiesseismicwavefrontgeometrylaw, def:physics:phyx-mini-0326:phyxminiproblems-problemphyxmini0326-satisfieslosslessseismicintensitylaw}
  Combining the energy, duration, cylindrical-area, and lossless-propagation laws yields the general surface-wave intensity formula at the station. No displayed answer is mentioned in this intermediate statement
\end{lemma}
\begin{proof}
  The conclusion follows from the governing relations and figure readouts named in the statement of surface Wave Intensity At Station formula.
\end{proof}
\begin{definition}[Answer Choice]
  \label{def:physics:phyx-mini-0326:phyxminiproblems-problemphyxmini0326-answerchoice}
  \lean{PhyXMiniProblems.ProblemPhyXMini0326.AnswerChoice}
  Labels of the four answers printed with the question
\end{definition}
\begin{proof}
  This is the declared model definition of Answer Choice.
\end{proof}
\begin{definition}[displayed Intensity In Kilowatts Per Square Meter]
  \label{def:physics:phyx-mini-0326:phyxminiproblems-problemphyxmini0326-displayedintensityinkilowattspersquaremeter}
  \lean{PhyXMiniProblems.ProblemPhyXMini0326.displayedIntensityInKilowattsPerSquareMeter}
  \uses{def:physics:phyx-mini-0326:phyxminiproblems-problemphyxmini0326-answerchoice}
  Displayed surface-wave intensities, in kilowatts per square metre
\end{definition}
\begin{proof}
  This is the declared model definition of displayed Intensity In Kilowatts Per Square Meter.
\end{proof}
\begin{definition}[recorded Dataset Answer Choice]
  \label{def:physics:phyx-mini-0326:phyxminiproblems-problemphyxmini0326-recordeddatasetanswerchoice}
  \lean{PhyXMiniProblems.ProblemPhyXMini0326.recordedDatasetAnswerChoice}
  \uses{def:physics:phyx-mini-0326:phyxminiproblems-problemphyxmini0326-answerchoice}
  The answer label recorded in the supplied dataset metadata
\end{definition}
\begin{proof}
  This is the declared model definition of recorded Dataset Answer Choice.
\end{proof}
\begin{definition}[Rounds To Displayed Whole Kilowatt Per Square Meter]
  \label{def:physics:phyx-mini-0326:phyxminiproblems-problemphyxmini0326-roundstodisplayedwholekilowattpersquaremeter}
  \lean{PhyXMiniProblems.ProblemPhyXMini0326.RoundsToDisplayedWholeKilowattPerSquareMeter}
  Agreement with a displayed whole-number intensity to the nearest unit
\end{definition}
\begin{proof}
  This is the declared model definition of Rounds To Displayed Whole Kilowatt Per Square Meter.
\end{proof}
% --- Archon declaration topology end ---
% --- Archon physics formalization source end ---
